\begin{titlepage}
\titresection{Résumé}

\vspace{0.75cm}
\begin{spacing}{1}
Le présent rapport analyse Raw Material Insight, un système d’aide à la décision développé au sein de la R\& D Pierre Fabre pour la gestion stratégique des matières premières. Face à la complexité croissante des exigences réglementaires, à l’augmentation continue du nombre de matières premières utilisées dans les formules et au besoin de sécuriser les processus industriels, l’entreprise a mis en place un système multi-source permettant de consolider les données du PLM, des spinners, des statistiques de ventes (SDV), de la Data Marketplace et d’outils internes tels que SharePoint.
Raw Material Insight fournit une vision synthétique et hiérarchisée du portefeuille matières, identifie les dossiers nécessitant une revue, calcule l’impact économique de chaque matière selon son utilisation en production, et met à disposition des outils de recherche avancés et des fiches techniques consolidées. Construit autour d’un flux ETL Alteryx automatisé et d’une couche décisionnelle BI, il est devenu un point d’entrée central pour différents métiers : homologation, sourcing, analytique, sécurisation et formulation.
Mon rôle dans l’étude du système a consisté en une observation approfondie de son fonctionnement et une compréhension de l’architecture afin de contribuer à la future extension analytique du projet, notamment l’intégration de JMP et la mise en place d’une base de données dédiée aux analytiques.
L’étude présentée dans ce rapport s’inscrit dans les attendus du génie logiciel : analyse du besoin, modélisation des données, étude des flux de production et présentation des modalités de mise en production.


\vspace{0.5em}
\end{spacing}

\vspace{2em}

\titresection{Abstract}


\vspace{0.75cm}
\begin{spacing}{1}
This report provides an analysis of Raw Material Insight, a decision-support system developed within Pierre Fabre’s R\&D to enhance the strategic management of raw materials. As regulatory demands increase and the number of raw materials grows, the company needed a multi-source digital system capable of consolidating PLM data, spinners, sales statistics (SDV), the Data Marketplace and internal repositories such as SharePoint.
Raw Material Insight offers a synthetic and prioritized view of the raw material portfolio, identifies dossiers requiring regulatory review, evaluates the economic impact of each material based on production usage, and provides advanced search features and consolidated technical sheets. Built on an automated Alteryx ETL workflow and a BI dashboard layer, the system has become a central tool across several R\& D functions, including homologation, sourcing, analytics, safety and formulation.
My contribution to the project focuses on the observation of the system, the exploration of its data architecture and the preparation for future analytical extensions, particularly through JMP and the construction of an analytical database.
This report follows software engineering methodology: needs analysis, data modeling, production workflow description and implementation details.

\end{spacing}
\end{titlepage}